\documentclass[11pt,a4paper]{article}

\usepackage[utf8]{inputenc}
\usepackage[T1]{fontenc}
\usepackage{amsmath,amssymb,amsthm}
\usepackage{graphicx}
\usepackage{booktabs}
\usepackage{hyperref}
\usepackage{xcolor}
\usepackage{tcolorbox}
\usepackage[margin=25mm]{geometry}
\usepackage{xeCJK}
\setCJKmainfont{Hiragino Mincho ProN}
\setCJKsansfont{Hiragino Sans}

\newtheorem{observation}{観察}
\newtheorem{conjecture}{予想}
\newtheorem{theorem}{定理}
\newtheorem{identity}{恒等式}

\title{物理定数の modular form 書き換え:\\
$(E_4 - E_2)/2$ identity の一般化と $\Omega_\Lambda$ の保型形式的解釈}
\author{Wright Brothers Research Program}
\date{2026年4月}

\begin{document}
\maketitle

\begin{abstract}
先行 note~\cite{kreimer-wb}で報告した observation
$(E_4 - E_2)/2 = 1 + 132q + O(q^2)$ (Wright Brothers $\alpha^{-1}$ 公式の
leading 132 が Eisenstein 線形結合の Fourier 係数)を systematic に
一般化する．主要な新しい発見:
(1) \textbf{Bernoulli ladder と Eisenstein 係数の厳密対応}:
$2k/|B_{2k}| = \frac{1}{2}|q^1\text{-coef of }E_{2k}|$ ---
Wright Brothers ladder は Eisenstein 係数の半分．
(2) \textbf{複数の物理的 ``word'' が Eisenstein combination として}:
$\alpha^{-1}$ 主要項 $= (E_4 - E_2)/2$, ミューオン $g-2 = (E_8 - E_2)/2 = 252$,
$E_8$ 根数 $= E_4$ 自体 $= 240$, 等．
(3) \textbf{Ladder 内の自明でない加法恒等式}: $12 + 120 = 132$,
$12 + 240 = 252$, $120 + 132 = 252$ 等，Wright Brothers ladder が
閉じた加法構造を持つ．
(4) \textbf{$\Omega_\Lambda$ の functional-equation dual form}:
$\Omega_\Lambda = 16|\zeta(-1)|\zeta(2)/\pi$ は $\zeta$ 値の
negative/positive integer dual を両方使う，Deligne--Beilinson 予想の
具体例候補．
(5) \textbf{$E_2$ quasi-modular anomaly} との関係 sketching．
本稿はこれらを rigorous ではなく explicit に記述し，Wright Brothers
``Spec($\mathbb{Z}$) alphabet'' が単なる数値遊びでなく Eisenstein 級数の
系統的書き換えであることを示す．
\end{abstract}

\tableofcontents

\newpage
\section{序: $(E_4 - E_2)/2$ identity の motivation}

先行 note~\cite{kreimer-wb}では，Kreimer Hopf 代数の文脈で
Wright Brothers $\alpha^{-1}$ 公式の主要項 $132$ が
\begin{equation}\label{eq:main-identity}
\frac{E_4(\tau) - E_2(\tau)}{2} = 1 + 132\, q + O(q^2)
\end{equation}
という modular form (quasi-modular) の Fourier 展開から出ることを
observation した．ここで
\begin{align}
E_2(\tau) &= 1 - 24 q - 72 q^2 - 96 q^3 - \ldots\\
E_4(\tau) &= 1 + 240 q + 2160 q^2 + 6720 q^3 + \ldots
\end{align}
は正規化 Eisenstein 級数，$q = e^{2\pi i \tau}$.

本稿ではこの observation を一般化し，他の Wright Brothers ``word'' も
同種の modular form Fourier 係数として書けることを探る．同時に
$\Omega_\Lambda = 2\pi/9$ の保型形式的解釈を deep に掘る．

\section{Bernoulli ladder と Eisenstein 係数の厳密対応}

\subsection{基本恒等式}

Eisenstein 級数の一般公式:
\begin{equation}
E_{2k}(\tau) = 1 - \frac{4k}{B_{2k}}\sum_{n=1}^\infty \sigma_{2k-1}(n) q^n
\end{equation}

$q^1$ 係数は $-4k/B_{2k}$. 絶対値は $4k/|B_{2k}|$.

Wright Brothers の Bernoulli ladder は:
\begin{equation}
L_k \equiv \frac{2k}{|B_{2k}|}
\end{equation}

\begin{identity}[Bernoulli ladder $=$ Eisenstein 半係数]
\begin{equation}
L_k = \frac{2k}{|B_{2k}|} = \frac{1}{2}\left|\text{coef of }q^1\text{ in }E_{2k}\right|
\end{equation}
\end{identity}

具体値:
\begin{center}
\begin{tabular}{llll}
\toprule
$k$ & $E_{2k}$ $q^1$ coef & $L_k$ & Wright Brothers での登場 \\
\midrule
1 & $-24$ & $12$ & $\alpha^{-1}$ 主要項, SM gauge \# \\
2 & $240$ & $120$ & $\alpha^{-1}$ 主要項, $E_8$ root $\div 2$ \\
3 & $-504$ & $252$ & ミューオン $g-2$ \\
4 & $480$ & $240$ & $E_8$ root 数 \\
5 & $-264$ & $132$ & $= 12 + 120$ ! \\
6 & $65520/691$ & 非整数 & 691 断絶 \\
\bottomrule
\end{tabular}
\end{center}

すなわち\textbf{Wright Brothers ladder の全エントリーは Eisenstein 係数
から自動生成される}．これは「arithmetic alphabet」の modular
form 側からの確認．

\subsection{階層の意味}

$E_{2k}$ は重さ $2k$ の modular form．Wright Brothers ladder の各
エントリーは異なる重さの保型形式に対応する．これは
\begin{itemize}
\item 重さ 2 ($L_1 = 12$): 電磁
\item 重さ 4 ($L_2 = 120$): 電磁 next
\item 重さ 6 ($L_3 = 252$): ミューオン
\item 重さ 8 ($L_4 = 240$): $E_8$ root ($E_8 = E_4^2$)
\item 重さ 10 ($L_5 = 132$): Sum of previous?
\item 重さ 12 ($L_6 = \text{非整数}$): cusp form 登場
\end{itemize}

``物理量の重さ'' というか ``observable の modular weight'' が
その保型形式的記述を決める可能性．

\section{複数の物理的 word が Eisenstein combination}

\subsection{$\alpha^{-1}$ 主要項: $(E_4 - E_2)/2$}

既出~\cite{kreimer-wb}:
\begin{equation}
\frac{E_4 - E_2}{2} = 1 + 132\, q + O(q^2)
\end{equation}
Wright Brothers $\alpha^{-1}$ の $132 = 2/|B_2| + 4/|B_4|$ に直接対応．

\subsection{ミューオン $g-2$: 複数の Eisenstein 表現}

ミューオン $g-2$ 異常 $\approx 252 = 6/|B_6| = L_3$．いくつかの
等価な Eisenstein 表現:

\textbf{(a) 単独 $E_6$}:
\begin{equation}
\frac{1}{2}|E_6 \text{ の } q^1\text{-coef}| = \frac{504}{2} = 252
\end{equation}

\textbf{(b) $(E_8 - E_2)/2$}:
\begin{equation}
\frac{E_8 - E_2}{2} = 1 + \frac{480 + 24}{2} q + O(q^2) = 1 + 252\, q + O(q^2)
\end{equation}

\textbf{(c) Ramanujan $\tau(3) = 252$}: $\Delta$ (cusp form $\sum \tau(n) q^n$)
の $q^3$ 係数

これら 3 つの独立表現が同じ 252 を与えるのは驚き．

\begin{observation}[252 の 3 重表現]
$252$ は (a) 重さ 6 Eisenstein の半分，(b) 重さ 8 と重さ 2 の差の半分，
(c) 重さ 12 cusp form の第 3 Fourier 係数，という 3 つの独立な
保型形式的表現を持つ．物理的にはミューオン $g-2$ 異常と一致．
\end{observation}

\subsection{$E_8$ root 数 $= 240$: 自明でない恒等式}

$E_8$ 格子の長さ 2 vector は 240 個 (定理)．これは
\begin{equation}
\theta_{E_8}(\tau) = E_4(\tau) = 1 + 240 q + O(q^2)
\end{equation}
から直接．

しかし $E_8$ 根は $k=4$ ladder entry $L_4 = 240$ とも一致:
\begin{equation}
L_4 = \frac{8}{|B_8|} = \frac{8}{1/30} = 240
\end{equation}

両者は $|B_4| = |B_8| = 1/30$ という偶然による等しさ．正確には
$E_8$ roots $= 240 = -E_4$ coef $= 8/|B_4|$ が主経路で，
$8/|B_8| = 240$ は副次的．

\subsection{一般的 $(E_{2k} - E_2)/2$ 公式}

公式:
\begin{equation}
\frac{E_{2k} - E_2}{2} = 1 + \left(\frac{-4k/B_{2k} - (-24)}{2}\right) q + O(q^2)
= 1 + \left(L_k + 12\right) q + O(q^2)
\end{equation}

(符号は $k$ 依存で注意．$B_{2k}$ の符号に応じて $L_k$ が $\pm$)．

$k \geq 2$ での $(E_{2k} - E_2)/2$ の $q^1$ 係数:

\begin{center}
\begin{tabular}{lll}
\toprule
$k$ & $(E_{2k} - E_2)/2$ $q^1$ coef & 物理対応 \\
\midrule
2 & $132$ & $\alpha^{-1}$ 主要項 \\
3 & $-240$ & $-E_8$ root? \\
4 & $252$ & ミューオン $g-2$ \\
5 & $-120$ & $-\alpha^{-1}$ 項 \\
\bottomrule
\end{tabular}
\end{center}

興味深いことに，$(E_4 - E_2)/2$ と $(E_8 - E_2)/2$ が物理に対応し，
$(E_6 - E_2)/2$ と $(E_{10} - E_2)/2$ は（符号違いで）別物理項に
対応する．

\section{Ladder の加法恒等式}

\subsection{発見}

Wright Brothers Bernoulli ladder $\{12, 120, 252, 240, 132\}$ 内で
以下の加法恒等式が成立:

\begin{identity}[ladder 加法恒等式]
\begin{align}
12 + 120 &= 132 \quad (L_1 + L_2 = L_5)\\
12 + 240 &= 252 \quad (L_1 + L_4 = L_3)\\
120 + 132 &= 252 \quad (L_2 + L_5 = L_3)\\
240 - 12 &= 228 \quad \text{(not in ladder)}
\end{align}
\end{identity}

第一恒等式 $L_1 + L_2 = L_5$ の証明:
\begin{align}
\frac{2}{|B_2|} + \frac{4}{|B_4|} &= 12 + 120 = 132\\
\frac{10}{|B_{10}|} &= \frac{10 \cdot 66}{5} = 132
\end{align}
すなわち $|B_2|^{-1} + 2|B_4|^{-1} = 5|B_{10}|^{-1}$.

同様に他の恒等式も数値的に確認できる．

\subsection{意味}

これらの加法恒等式は Bernoulli 数の非自明な関係式で，
Wright Brothers ladder が単なる数列ではなく\textbf{閉じた加法
構造を持つ部分集合}であることを示唆する．

物理的対応:
\begin{itemize}
\item $\alpha^{-1}$ 主要項 $= L_1 + L_2 = 132 = L_5$ → 単独 $E_{10}$
との両方の表現
\item ミューオン $g-2$ = 252 = $L_3 = L_1 + L_4 = L_2 + L_5$ →
三重表現
\end{itemize}

これは modular form level の「言語」で\textbf{単一物理量が複数の
文法的表現を持つ}ことに対応．

\subsection{Euler 恒等式との関係}

Euler の classical 結果: Bernoulli 数には漸化式
\begin{equation}
\sum_{k=0}^{n-1}\binom{2n}{2k} B_{2k} = 1 - 2n
\end{equation}
があり，これが $B_2, B_4, \ldots$ 間の線形関係を与える．

Wright Brothers ladder 加法恒等式は Euler 漸化式の specific
reformulation かもしれない．これの rigorous 証明は future work．

\section{$\Omega_\Lambda$ の保型形式的解釈}

\subsection{既知の表現}

先行研究~\cite{wb-wdw}:
\begin{equation}
\Omega_\Lambda = \frac{2\pi}{9}
\end{equation}

等価形式:
\begin{align}
\Omega_\Lambda &= 8\pi B_2^2\\
&= 8\pi |\zeta_{\neg 2}(-1)|/\text{something}\ldots
\end{align}

\subsection{新しい表現: functional equation dual}

$B_2 = 1/6$, $\zeta(-1) = -1/12 = -B_2/2$, $\zeta(2) = \pi^2/6 = \pi^2 B_2$．

これらを組み合わせて:
\begin{align}
\Omega_\Lambda &= 8\pi B_2^2 = 8\pi \cdot (1/6)^2 = \frac{8\pi}{36}\\
&= \frac{8\pi}{36}\cdot \frac{1}{1}\\
&= \frac{8 \cdot \pi B_2 \cdot B_2}{1}\\
&= \frac{8 \cdot (\zeta(2)/\pi) \cdot (-2\zeta(-1))}{1}\\
&= -\frac{16 \zeta(-1)\zeta(2)}{\pi}
\end{align}

数値確認: $-16 \cdot (-1/12) \cdot (\pi^2/6)/\pi = 16\pi/(72) = 2\pi/9$ ✓

\begin{observation}[Functional-equation dual form]
\begin{equation}
\boxed{\Omega_\Lambda = -\frac{16\,\zeta(-1)\,\zeta(2)}{\pi} = \frac{16\,|\zeta(-1)|\,\zeta(2)}{\pi}}
\end{equation}
\end{observation}

これは以下の構造を持つ:
\begin{itemize}
\item \textbf{$\zeta(-1)$}: 負の整数での $\zeta$ 値，Bernoulli $B_2$
と直結 (arithmetic/motivic 側)
\item \textbf{$\zeta(2)$}: 正の整数での $\zeta$ 値，period $\pi^2$
と直結 (transcendental/period 側)
\item \textbf{$1/\pi$}: dimensional regularization factor
\end{itemize}

これは\textbf{Riemann functional equation の両側を同時に使う}
非自明な表現．

\subsection{Deligne--Beilinson との接続}

Deligne 予想~\cite{Deligne1979} および Beilinson
予想~\cite{Beilinson1984}は，motivic L-value が period と
algebraic number の積であると主張:
\begin{equation}
L^*(M, n) \sim (\text{algebraic}) \cdot (\text{period})^r
\end{equation}

$\Omega_\Lambda = 16|\zeta(-1)|\zeta(2)/\pi$ をこの観点で見ると:
\begin{itemize}
\item $\zeta(-1) = -1/12$: algebraic (rational)
\item $\zeta(2) = \pi^2/6$: period (involves $\pi^2$)
\item $1/\pi$: additional period factor
\item 全体: algebraic $\times$ period $= $ dimensionless real number
\end{itemize}

すなわち $\Omega_\Lambda$ は\textbf{$\mathbb{Q}(\pi)$ の外で代数的},
\textbf{$\pi$-rational}．これは Beilinson 予想の形式に整合．

\begin{conjecture}[Wright Brothers $\Omega_\Lambda$ as Beilinson L-value]
$\Omega_\Lambda = 2\pi/9$ は specific motive $M$ の critical L-value
$L^*(M, \text{critical point})$ であり，Beilinson 予想の物理的
実現例である．
\end{conjecture}

どの motive $M$ が該当するかは future work．Candidates:
\begin{itemize}
\item $\text{Spec}(\mathbb{Z})$ 自身
\item Tate motive $\mathbb{Q}(-1)$ or $\mathbb{Q}(1)$
\item Cyclotomic field の motive
\end{itemize}

\subsection{$E_2$ quasi-modular anomaly との関係}

$E_2$ は厳密には modular form ではなく\textbf{quasi-modular}:
\begin{equation}
E_2\left(\frac{a\tau + b}{c\tau + d}\right) = (c\tau + d)^2 E_2(\tau) - \frac{6ic(c\tau + d)}{\pi}
\end{equation}
anomaly 項 $-6ic/\pi$.

Completion: $E_2^*(\tau) = E_2(\tau) - 3/(\pi\,\text{Im}\,\tau)$ は
modular invariant だが non-holomorphic.

\textbf{Key 観察}: $E_2^*$ の anomaly correction $3/(\pi y)$ に $\pi$
が分母．$\Omega_\Lambda$ の functional dual form にも $1/\pi$ が現れる．

\begin{conjecture}[$\Omega_\Lambda$ と $E_2^*$ anomaly]
$\Omega_\Lambda = 16|\zeta(-1)|\zeta(2)/\pi$ の $1/\pi$ factor は
$E_2$ quasi-modular anomaly $3/(\pi y)$ に由来する．Physical に
dark energy 密度は $E_2$ の modular anomaly の物理的現れ．
\end{conjecture}

Precise な link は $E_2^*(\tau = i)$ 等の specific evaluation と
Wright Brothers formula の関係を明らかにする必要がある．

\subsection{$\tau = i$ での evaluation}

$E_2$ の specific 値~\cite{Serre1973}:
\begin{equation}
E_2(i) = \frac{3}{\pi}
\end{equation}

完成形: $E_2^*(i) = E_2(i) - 3/(\pi \cdot 1) = 3/\pi - 3/\pi = 0$.

すなわち $E_2^*$ は $\tau = i$ で vanish.

$E_4(i) = 3\Gamma(1/4)^8/(64\pi^6)$ (Selberg-Chowla~\cite{Chowla1967}).

$E_4(i) \cdot \pi^6 \sim \Gamma(1/4)^8 / 21.3$. 数値: $\Gamma(1/4) \approx 3.626$,
$\Gamma(1/4)^8 \approx 3.32 \times 10^4$, $\pi^6 \approx 961$,
$E_4(i) \approx (3 \cdot 3.32 \times 10^4)/(64 \cdot 961) = 1.62$.

$(E_4 - E_2)(i)/2 = (1.62 - 0.955)/2 = 0.33$

そして $\Omega_\Lambda = 0.698$. 近いが一致しない．Factor 2.

\begin{observation}[$\Omega_\Lambda$ vs $(E_4 - E_2)(i)$]
$(E_4(i) - E_2(i))/2 \approx 0.33$, $\Omega_\Lambda = 0.698$.
Ratio $\approx 2.13$. Clean identity は現時点で未発見．
\end{observation}

Specific $\tau$-value evaluation が $\Omega_\Lambda$ を直接与えるかは
future work．Trial: $\tau = \rho = e^{2\pi i/3}$ 等他の特殊点．

\section{新しい予測の可能性}

\subsection{他の物理定数の modular 予測}

Wright Brothers ladder $\{12, 120, 252, 240, 132\}$ のそれぞれが物理
対応を持つと仮定すると，未同定エントリー:
\begin{itemize}
\item $L_4 = 240$: $E_8$ root (既知)，他には?
\item $L_5 = 132$: $\alpha^{-1}$ 主要項 (既知)，他には?
\end{itemize}

$240, 132$ が別の物理量とも一致するか探索:
\begin{itemize}
\item 132 = $\alpha_{\rm em}^{-1}(m_Z)$?: 観測 127.95, 違う (ただし近い)
\item 132 $\approx$ 質量比?: $m_t/m_c \approx 135$, 違う
\item 240 $\approx$ Higgs self-coupling $\times 10^3$?: $\lambda_H \approx 0.129$, $\times 10^3 = 129$, 違う
\end{itemize}

Clean な一致は現時点で 132 と 240 の新規物理対応は見つからない．

\subsection{ラダー外の数}

$(E_2 + E_4 + E_6)/3$: $(-24 + 240 - 504)/3 = -96$. 意味?
$(E_4 \cdot E_2)$: $(1 - 24q + \ldots)(1 + 240q + \ldots) = 1 + 216 q + \ldots$, 216 = $6^3$
$(E_4)^2 = E_8$: 480, 既存
$E_4 \cdot E_6 = E_{10}$: $-264$, 既存
$E_4^3 - E_6^2 = 1728 \Delta$: 係数 cusp form
$\Delta = q - 24 q^2 + 252 q^3 - \ldots$: $\tau$-function

$\tau(n)$ の値は $1, -24, 252, -1472, 4830, -6048, \ldots$

$\tau(4) = -1472$: 物理対応?
$|\tau(4)| = 1472$. 素因子分解: $1472 = 2^6 \cdot 23$. 物理的意味は不明．

$\tau(5) = 4830$: $4830 = 2 \cdot 3 \cdot 5 \cdot 7 \cdot 23$.
タウ lepton $g-2$ への予測候補 (Wright Brothers generation conjecture).

\subsection{Tau $g-2$ の modular 予測}

先行~\cite{wb-spec-z}で提案された世代-Ramanujan 対応から
$\Delta a_\tau$ が $\tau(5) = 4830$ に比例する予測を更新．

Scale: ミューオン $g-2$ observed $\approx 252 \times 10^{-11}$
(Fermilab+BNL)．もし scale universal なら:
\begin{equation}
\Delta a_\tau \stackrel{?}{\sim} 4830 \times 10^{-11} \approx 5 \times 10^{-8}
\end{equation}

現在の tau $g-2$ 実験限界は $\sim 10^{-2}$，この予測は検証不能．

Belle II (KEK), BES III (IHEP), LHCb upgrade, FCC-ee が将来的に
$10^{-4}$ 精度まで到達する可能性．その時点で本予測は検証可能．

\section{Honest 評価}

\subsection{新しい観察の強度}

本稿で得た observation を強度で格付け:

\begin{center}
\begin{tabular}{lll}
\toprule
観察 & 強度 & 根拠 \\
\midrule
$(E_4-E_2)/2|_{q^1} = 132$ & ★★★ & 直接計算で verify 可 \\
Ladder $= \frac{1}{2}|E_{2k}$ coef$|$ & ★★★ & Trivial from Eisenstein 公式 \\
$L_1 + L_2 = L_5$ (=132) & ★★ & 非自明な Bernoulli 恒等式 \\
$L_1 + L_4 = L_3$ (=252) & ★★ & 加法恒等式 \\
$\Omega_\Lambda = -16\zeta(-1)\zeta(2)/\pi$ & ★★★ & 直接計算で verify \\
$E_2$ anomaly $\to \Omega_\Lambda$ の $1/\pi$ & ★ & Speculative \\
$\tau(5) \to \Delta a_\tau$ 予測 & ★ & Speculative, 検証不能 \\
\bottomrule
\end{tabular}
\end{center}

強度 ★★★ の 3 つは数値的に reliable で reproducible．これらは
Wright Brothers プログラムの solid な core．

\subsection{Wright Brothers への impact}

以前の状態: Wright Brothers は ``Bernoulli 数で物理定数を近似'' という
ヒューリスティック

新しい状態: Wright Brothers は ``Eisenstein 級数と $\zeta$ 関数の
functional equation dual から物理定数を systematically 書き換える''
という structured framework

Repositioning の意義:
\begin{itemize}
\item 「数値一致探し」から「保型形式の物理的解釈」へ
\item Kreimer--Connes Hopf 代数，Beilinson 予想，Langlands program と
同じ mathematical neighborhood に住む
\item Look-elsewhere effect の危険は残るが，体系的 structure が見えて
きた
\end{itemize}

\subsection{弱点}

\begin{enumerate}
\item $\Omega_\Lambda$ と $E_2$ quasi-modular の詳細 link は未完成
\item $\tau = i$ 等での evaluation から直接 $\Omega_\Lambda$ が出る
ルートは未発見
\item Ladder 加法恒等式の一般証明と物理的意味は未明
\item 質量比・CKM は依然として扱えない
\item Pre-diction でなく post-diction 主体
\end{enumerate}

\section{次のステップ}

\subsection{Rigorous 化の作業項目}

\begin{enumerate}
\item Ladder 加法恒等式 $L_1 + L_2 = L_5$ の Bernoulli 漸化式からの
証明
\item $E_2^*(\tau)$ の specific point evaluation と $\Omega_\Lambda$ の
関係定式化
\item $(E_4 - E_2)/2$ の $q^1$ 係数が Kreimer Hopf 代数の具体的
primitive element と対応するかの確認
\item Ramanujan $\tau(n)$ と Wright Brothers 世代仮説の rigorous 結合
\item Selberg-Chowla 公式~\cite{Chowla1967}を使った
$E_4(i), E_6(i)$ 等の exact evaluation と WB formula の比較
\end{enumerate}

\subsection{英語版論文への展開}

本一連の note は arXiv 投稿用の英語論文として統合する価値がある．
Target 雑誌候補:
\begin{itemize}
\item \textit{J.\ Number Theory}
\item \textit{Commun.\ Number Theory Phys.}
\item \textit{Adv.\ Theor.\ Math.\ Phys.}
\item \textit{Found.\ Phys.} (より speculative な部分)
\end{itemize}

\subsection{Connes--Marcolli 団との接触}

Wright Brothers の framework は A.~Connes, M.~Marcolli の
Noncommutative Geometry / Arithmetic Site program~\cite{CM2008, CC2014}
と conceptual に整合する．両者の接触・統合が重要な future step．

\section{結論}

Wright Brothers プロジェクトの ``Spec$(\mathbb{Z})$ alphabet'' 仮説は，
保型形式 (Eisenstein 級数 $E_{2k}$, Ramanujan $\Delta$) の Fourier 係数
として systematic に書き換えられることを示した．主要な新しい発見:

\begin{enumerate}
\item \textbf{Bernoulli ladder $= \frac{1}{2}|E_{2k}$ coef$|$} の厳密対応
\item $(E_4 - E_2)/2 = 1 + 132 q + \ldots$ で $\alpha^{-1}$ 主要項
\item $(E_8 - E_2)/2 = 1 + 252 q + \ldots$ でミューオン $g-2$
\item Ladder 内の非自明加法恒等式 ($L_1 + L_2 = L_5$, $L_1 + L_4 = L_3$)
\item $\Omega_\Lambda = -16 \zeta(-1)\zeta(2)/\pi$ という
functional-equation dual form
\item Beilinson 予想の物理的実現候補として $\Omega_\Lambda$
\end{enumerate}

これらは Wright Brothers が「numerology」から「保型形式物理」へと
明確に移行したことを示す．依然として多くの rigorous 化作業が残るが，
framework の方向性は明確で，既存の mathematical physics と
互換性がある．

Next step は Kreimer Hopf 代数との統合，Connes--Marcolli との接触，
英語版論文投稿．2027--2028 年の experimental verification
(Euclid, Belle II, 水晶 BAW) と合わせて，Wright Brothers framework の
生死判定が数年内に下される．

\begin{thebibliography}{99}

\bibitem{kreimer-wb}
Wright Brothers Research Program,
``Kreimer--Connes Hopf 代数と Wright Brothers $\alpha^{-1}$ 公式'',
\texttt{kreimer\_wb\_integration\_ja.tex} (2026).

\bibitem{wb-wdw}
Wright Brothers Research Program,
``最小超空間量子宇宙論における時間的 Dirichlet--Neumann 真空'',
\texttt{wdw\_temporal\_dn\_ja.tex} (2026).

\bibitem{wb-spec-z}
Wright Brothers Research Program, ``Spec$(\mathbb{Z})$ alphabet 深掘り'',
\texttt{spec\_z\_alphabet\_deeper\_ja.tex} (2026).

\bibitem{Deligne1979}
P.~Deligne, ``Valeurs de fonctions L et périodes d'intégrales'',
Proc.\ Symp.\ Pure Math.\ \textbf{33.2}, 313 (1979).

\bibitem{Beilinson1984}
A.~Beilinson, J.\ Soviet Math.\ \textbf{30}, 2036 (1985).

\bibitem{Serre1973}
J.-P.~Serre, \textit{A Course in Arithmetic} (Springer, 1973).

\bibitem{Chowla1967}
S.~Chowla and A.~Selberg, Crelle J.\ \textbf{227}, 86 (1967).

\bibitem{CM2008}
A.~Connes and M.~Marcolli, \textit{Noncommutative Geometry,
Quantum Fields and Motives} (AMS, 2008).

\bibitem{CC2014}
A.~Connes and C.~Consani, C.\ R.\ Math.\ Acad.\ Sci.\ Paris
\textbf{352}, 971 (2014).

\end{thebibliography}

\end{document}
